\documentclass{beamer}
\usepackage[spanish]{babel}
\usepackage[T1]{fontenc}
\usepackage[utf8]{inputenc}
\usepackage{lmodern}
\usepackage{amssymb}
\usepackage{graphicx}
\usepackage{amsthm}
\usepackage{amsmath}
\usepackage{mathtools}


\uselanguage{spanish}
\languagepath{spanish}


\def\tituloTesis{Métricas de mimetización acústico-prosódica en hablantes y su relación con rasgos sociales de diálogos}

\title{\tituloTesis}
\author{Juan Manuel Pérez}

\usetheme{Madrid}

\begin{document}

\beamertemplatenavigationsymbolsempty


\frame{\titlepage}



\begin{frame}
  \frametitle{¿Cómo dijo?}


\begin{itemize}
  \item Sistemas de diálogo humano-computadora son cada vez más frecuentes, y sus aplicaciones comprenden una amplia gama de rubros
  \item ``Bien'' en la dimensión lingüística, mal en todo lo suprasegmental: emociones, actitudes, intenciones.
  \item Mimetización: Fenómeno insconsciente que se manifiesta a través de la adaptación de los hablantes. Fuertemente emparentada con el sentimiento de empatía.
  \item Objetivo del trabajo: Explorar y refinar una métrica de la mimetización acústico-prosódica, y validar que capture ciertas percepciones sociales en un corpus de diálogos en inglés.
\end{itemize}

\end{frame}



\section{Introducción}

\begin{frame}
  \frametitle{Sistemas de diálogo Humano-Computadora}

  \begin{enumerate}[<+->]
    \item Sistemas actuales
    \item Bien en la parte lingüística de la comunicación: entender y transmitir mensajes estructuralmente correctos.
    \item Mal en la parte superestructural: intercambio de emociones, actitudes, etc.
    \item El presente trabajo trata de hacer un (pequeño) aporte sobre el análisis de la ``naturalidad'' de las conversaciones.
  \end{enumerate}
\end{frame}


\begin{frame}
  \begin{columns}
    \column{0.5\textwidth}
    \frametitle{Ejemplos de ``falta de naturalidad''}
    \begin{enumerate}
      \item Sistemas de llamadas comerciales
      \item Siri, Google Now
      \item Otros?
    \end{enumerate}
    \column{0.5\textwidth}
    \begin{figure}
      \includegraphics[scale=0.25]{images/hal.jpg}
    \end{figure}
  \end{columns}
\end{frame}


\begin{frame}
  \frametitle{Prosodia}
  \framesubtitle{Definir acá qué es la prosodia aproximadamente}

\end{frame}



\begin{frame}
  \frametitle{Entrainment}

  \begin{enumerate}
    \item Fenómeno ubícuo en la comunicación
    \item Ocurre a varios niveles
    \item Largamente estudiado en psicología de comportamiento (referencias)
    \item (ya está, la próxima conversación que tengan afuera de acá van a chequear ésto)
  \end{enumerate}
\end{frame}

\begin{frame}
  \frametitle{¿Y cómo lo medimos?}
  La definición de entrainment hasta acá vista es heurística! ¿Cómo definimos una medida para esto?

  Vamos a explorar una métrica definida en trabajos anteriores, pulirla un poco, y verificar que efectivamente capture ciertas características del entrainment.

  ¿Cómo? Usando un corpus con anotaciones sociales


\end{frame}


\section{Antecedentes}

\begin{frame}
  \frametitle{Otras métricas}

  \begin{itemize}
    \item La mimetización es un fenómeno tanto lineal: se va acentuando a lo largo de la conversación
    \item Pero también es un fenómeno dinámico: va variando localmente a lo largo de la conversación.
  \end{itemize}

  Muchas métricas sólo toman la parte global, dividiendo la conversación en 2 o más partes y luego calculando la diferencia entre las medias de las diferentes variables acústicas en cada sección.
\end{frame}


\begin{frame}
  \frametitle{Problema del alineamiento de tiempo}

  \begin{figure}[t]
    \includegraphics[scale=0.40]{images/conversation_turns.pdf}
  \end{figure}

  \begin{itemize}
    \item ¿Cómo comparamos los diferentes turnos de una conversación?
    \item Comparar uno a uno es un enfoque simplista y no representativo de la realidad
  \end{itemize}
\end{frame}

\begin{frame}
  \frametitle{Método TAMA}
\end{frame}





\section{Materiales y Métodos}

\subsection{Corpus}
\begin{frame}
  \frametitle{Columbia Games Corpus}
  \framesubtitle{Descripción}
  \begin{columns}
    \column{0.35\textwidth}
      \begin{figure}
        \includegraphics[width=\textwidth]{images/columbia_games_color.jpg}
      \end{figure}

    \column{0.65\textwidth}

    \begin{itemize}
      \item Corpus de conversaciones diádicas en Inglés Americano Estándar
      \item 12 sesiones con 14 tareas/juegos cada una.
      \item En cada sesión, se sentó a dos participantes en una cabina profesional de grabación, y una cortina opaca colgando entre ellos para evitar la comunicación visual.
      \item Los participantes contaron con computadoras a través de las cuales interactuaban mediante juegos.
    \end{itemize}
  \end{columns}

\end{frame}


\begin{frame}
  \frametitle{Columbia Games Corpus}
  \framesubtitle{Juegos de objeto}
  \begin{columns}
    \column{0.35\textwidth}
      \begin{figure}
        \includegraphics[width=\textwidth]{images/columbia_games_color.jpg}
      \end{figure}

    \column{0.65\textwidth}

    \begin{itemize}
      \item Dos roles: Descriptor y Seguidor
      \item En la pantalla, se ven entre 5 y 7 objetos en posiciones aleatorias
      \item El Descriptor ve uno más, titilante, del cual debe describir su posición
      \item El Seguidor debe mover la representación del objeto a la misma posición de la pantalla
      \item Finalmente, se puntúa de 1 a 100 el posicionamiento del objeto
    \end{itemize}
  \end{columns}
\end{frame}


\begin{frame}
  \frametitle{Columbia Games Corpus}
  \framesubtitle{Detalles y Anotaciones sociales}

\end{frame}




\end{document}
